\documentclass{article}
\usepackage[utf8]{inputenc}
\usepackage[brazil]{babel}
\usepackage[a4paper, top=3cm, bottom=2cm, left=2cm, right=2cm]{geometry}

\title{aula 3 - respostas}
\author{marina}
\date{\today}

\begin{document}

\maketitle
\section{sistema auditivo}
tudo o que é do corpo e percebe sons. dividido em central e periférico.
\textbf{central:} percebe os sinais enviados pelo periférico e interpreta, associa esses sinais com símbolos
\textbf{periférico:} ouvido externo, médio e interno. capta o som, gera sinais eletroquímicos e envia pro cérebro.
\subsection{periférico}
\begin{itemize}
    \item orelha externa: o que é visto, captar e direcionar som, faz o som ambiente entrar no conduto auditivo
    \item conduto: protege a membrana do tímpano
    \item membrana timpanica: interfa
    \item orelha média: 3 ossículos, servem para fazer o casamento de impedância do ar com a água. som vê uma impedância maior na água do que no ar e tem uma dificuldade, ossículos faz isso ser mais fácil. também protegem o sistema auditivo quando somos expostos a sons muito elevados, músculos contraem e enrijegem para proteger
    \item janela oval: estribo faz parede da janela oval que vibra o líquido sinovial que fica na cóclea
    \item cóclea: escala média vibra que tem as células
    \item células ciliadas: abrem uma passagem para os íons do líquido sinovial, o íon entra no núceo da terra  eo valor vai para
\end{itemize}
\section{audibilidade}
audição não responde igual para todas as faixas de frequência \\
curvas isofônicas.\\
largura que responde a uma mesma frequência. se duas freq muitas poroxi,am, a cólcea na~cõ osnsgue i\\
2 sons próximos -> nao escitam sensações distintas, é difícil de separar\\
sons mais diferentes/distante -> consegue oucir melhor os dois sons\\
bandas críticas: largura de banda que o ouvido consegue separar\\
livro: brains\\
\section{exercício 01}
It must be pointed out that the measurements taken so far indicate that the critical bands have a certain width, but that their position on the frequency scale is not fixed; rather, the position can be changed continuously, perhaps by the ear itself.\\
A difference of 1 Bark corresponds to the width of one critical band over the whole frequency range, and also corresponds very nearly to a pitch interval of 100 mels.\\
mel: percepção entre altura e frequência de forma linear. mapeamento de percepção física e frequência\\
1.a) aumenta de forma não linear. banda levemente constante, até uma frequência média e depois aumenta\\
1.b) quanto maior a frequência, menos temos sensibilidade\\
1.c)

\end{document}